Educational satellites like the CubeSats reside almost exclusively at the \ac{LEO}, which ranges from 200~km to 2000~km above the earth surface, which has various reasons.
\newline
CubeSats are very small satellites and therefor the space for electronic devices and accumulators is very limited. Thus the short range between the device and \acp{GCC} is ideal as it reduces the power consumtion to a minimum and is close enough to earth so that the CubeSat will not be affected by the radiation of the Van Allen radiation belts \cite{horne_wave_2005}.
\newline
Moreover the deployment of a satellites in \ac{LEO} significantly decreases the costs compared to satellites stationed in the \ac{MEO} or \ac{GEO}. 
%da dann die preise rein
\newline
Furthermore, the closer the small satellite are to earth, the they will burn to ashes caused by entering earth's atmosphere. This helps keeping the orbits clean, as space debris can cause collateral damage to future space missions \cite{akahoshi_influence_2008}.
\newline
However, the small distance a satellite has in \ac{LEO} leads to a 